\section{Evaluation}
\label{sec:eval}

All results below were produced on the test laptop (12 logical CPUs, 15 GB RAM
) on 19--20 November 2026, using the scripts in the repository; the raw outputs are in \code{outputs/evaluation/}.

\subsection{Functional correctness}
\begin{center}\small
\begin{tabularx}{\linewidth}{@{}lXr@{}}
\toprule
Suite & What it exercises & Result \\
\midrule
Chaincode unit tests & every contract rule on an in-memory ledger that, like Fabric, gives a transaction no read-your-writes & 18/18 \\
Gateway unit tests & JS fingerprints byte-identical to the Python reference (three example manifests and a Unicode-heavy one); encrypted store (tamper, overwrite, address binding) & 10/10 \\
API end-to-end & the acceptance scenario of the roadmap on the live BFT network: registration of a photograph, a notebook and a knife; transfers between all three agencies; automatic identity checks; examination; court submission; verification; downloads; proofs; court outputs; independent verification of the bundle & 46/46 \\
Browser end-to-end & the same scenario through the UI in headless Chromium, as four different users, including role-dependent visibility of actions and printing of the court outputs to PDF & 26/26 \\
BFT network suite & consensus type and consenters; quorum signatures on blocks; tolerance of one stopped orderer; halt with two; recovery; leader failure and view change; persistence & 8/8 \\
\bottomrule
\end{tabularx}
\end{center}

\paragraph{Identity checks in the application.} The application reproduces the numbers of the worked examples (\cref{sec:worked}) because it calls the same reference tool: a full re-scan of the six notebook pages matched with a largest pHash distance of 4 (threshold 10); a re-scan with page 4 substituted was rejected, naming page 4; a re-scan with a page missing was rejected on page count; a 3D re-scan of the knife matched (RMS 0.568~mm, maximum 2.387~mm) and a re-scan with the tip broken off was rejected (maximum 5.024~mm).

\subsection{Block validation (D-030, D-050)}
Decoding the deployed channel configuration shows that peers accept a block only if it carries signatures of \textbf{3 of the 4 consenter identities} (\code{n\_out\_of: 3} over the four consenter certificates). Fabric 3.1.5's \code{configtxgen} writes this policy itself for BFT channels, replacing the \code{ANY Writers} rule of the sample configuration; generating the channel from the original configuration gives the same policy. Every observed block carries three consenter signatures. Peer \emph{rejection} of a forged single-signature block was not demonstrated, because producing one requires a modified orderer.

\subsection{Tamper suite}
\Cref{tab:tamper} lists attacks against stored data at rest and against the contract's inputs, run against the live network by \code{gateway-v2/scripts/tamper-suite.js}. All were detected or refused.
{\small
\begin{longtable}{@{}p{0.30\linewidth}p{0.24\linewidth}p{0.32\linewidth}l@{}}
\caption{Tamper suite: attack, expected behaviour, observed behaviour (gateway-v2/scripts/tamper-suite.js).}\label{tab:tamper}\\
\toprule
Attack & Expected & Observed & \\ \midrule\endhead
One bit of a stored (encrypted) evidence file flipped on disk & download refused & \footnotesize HTTP 409: evidence store integrity failure for A001: Unsupported state or unable to authenticate dat & \textbf{detected}\\
Stored manifest corrupted on disk & verification INCONCLUSIVE (never MATCH); store error reported & \footnotesize INCONCLUSIVE; evidence store integrity failure for the manifest of GJ-GNR-CYBPS-2026 & \textbf{detected}\\
Insider stores a doctored photo plus a self-consistent forged manifest & NO\_MATCH: the ledger still points to the original manifest & \footnotesize NO\_MATCH (ledger manifest 9a63dc08a1f8\ldots{}, forged fa02e40cb658\ldots{}) & \textbf{detected}\\
Valid agency certificate without a CustodyChain role writes to the ledger & rejected by the chaincode & \footnotesize caller User1@\allowbreak{}police.\allowbreak{}custodychain.\allowbreak{}local@\allowbreak{}PoliceMSP has no CustodyChain role & \textbf{detected}\\
Caller puts a different custodian, registrant and date in the arguments & ignored: actor from certificate, time from transaction & \footnotesize custodian io.\allowbreak{}sharma@\allowbreak{}PoliceMSP, registered\_at 2026-11-19T16:42:55.\allowbreak{}462Z & \textbf{detected}\\
Replay of the same registration & rejected (already exists) & \footnotesize evidence GJ-GNR-CYBPS-2026-775621-E900 already exists & \textbf{detected}\\
Receiver completes a Police-\textgreater{}FSL hand-over with FSL+Court endorsements only & rejected at validation (key-level policy) & \footnotesize Transaction be9d46820396\ldots{} failed to commit with s & \textbf{detected}\\
Same-agency colleague (not custodian) disposes of the item & rejected (custodian only) & \footnotesize only the current custodian may record DISPOSE & \textbf{detected}\\
File altered between the officer's browser and the gateway & rejected before anything is stored or anchored & \footnotesize HTTP 400: file 'IMG\_0001.\allowbreak{}jpg' changed in transit: browser SHA-256 2e2864e31f9b78 & \textbf{detected}\\
\bottomrule
\end{longtable}

}

\subsection{Performance: BFT versus a Raft control}
\label{sec:bench}
\paragraph{Method (D-051).} A load generator using the application's own Fabric client submits transactions as \code{io.sharma} at concurrency 1, 4, 8, 16 and 32 (at least 40 transactions per level) and measures end-to-end latency (endorsement by two peers, ordering, commit on the gateway's peer) and committed throughput. Three operations are measured: an access-log write (\code{RecordEvent(ACCESS)}, one key), a registration (\code{RegisterEvidence}: record, event, two indexes, key-level policy) and a query (\code{GetEvidence}). To isolate the cost of Byzantine fault tolerance, the identical network (4 orderers, 3 peers, same chaincode and host) was rebuilt with Raft (etcdraft, 4 consenters) and measured the same way, with batching aligned (50~ms / 100 requests).

\begin{table}[htbp]\centering\footnotesize\setlength{\tabcolsep}{4pt}
\begin{tabular}{@{}rrrrrrrrr@{}}
\toprule
 & \multicolumn{4}{c}{BFT, 4 orderers} & \multicolumn{4}{c}{Raft, 4 orderers (control)}\\
\cmidrule(lr){2-5} \cmidrule(lr){6-9}
Conc. & p50 ms & p95 ms & max ms & tx/s & p50 ms & p95 ms & max ms & tx/s\\
\midrule
1 & 64 & 122 & 122 & 14.19 & 105 & 121 & 124 & 9.59\\
4 & 70 & 127 & 135 & 46.1 & 117 & 141 & 146 & 33.68\\
8 & 129 & 153 & 157 & 60.84 & 107 & 133 & 135 & 72.98\\
16 & 128 & 215 & 253 & 116.46 & 135 & 157 & 172 & 118.23\\
32 & 149 & 209 & 276 & 198.35 & 162 & 204 & 232 & 195.61\\
\bottomrule
\end{tabular}
\caption{RecordEvent(ACCESS): latency and committed throughput by concurrency (failed transactions: bft4-r2 0, raft4-r2 0).}
\label{tab:bench-recordeventaccess}
\end{table}

\begin{table}[htbp]\centering\footnotesize\setlength{\tabcolsep}{4pt}
\begin{tabular}{@{}rrrrrrrrr@{}}
\toprule
 & \multicolumn{4}{c}{BFT, 4 orderers} & \multicolumn{4}{c}{Raft, 4 orderers (control)}\\
\cmidrule(lr){2-5} \cmidrule(lr){6-9}
Conc. & p50 ms & p95 ms & max ms & tx/s & p50 ms & p95 ms & max ms & tx/s\\
\midrule
1 & 65 & 74 & 106 & 15.46 & 105 & 118 & 123 & 9.63\\
4 & 73 & 136 & 137 & 45.51 & 119 & 133 & 135 & 35.27\\
8 & 116 & 163 & 171 & 63.16 & 125 & 166 & 173 & 62.97\\
16 & 122 & 161 & 188 & 130.73 & 151 & 231 & 301 & 99.96\\
32 & 172 & 213 & 281 & 182.11 & 221 & 366 & 427 & 131.45\\
\bottomrule
\end{tabular}
\caption{RegisterEvidence: latency and committed throughput by concurrency (failed transactions: bft4-r2 0, raft4-r2 0).}
\label{tab:bench-registerevidence}
\end{table}

\begin{table}[htbp]\centering\footnotesize\setlength{\tabcolsep}{4pt}
\begin{tabular}{@{}rrrrrrrrr@{}}
\toprule
 & \multicolumn{4}{c}{BFT, 4 orderers} & \multicolumn{4}{c}{Raft, 4 orderers (control)}\\
\cmidrule(lr){2-5} \cmidrule(lr){6-9}
Conc. & p50 ms & p95 ms & max ms & tx/s & p50 ms & p95 ms & max ms & tx/s\\
\midrule
1 & 6 & 8 & 10 & 163.32 & 11 & 13 & 14 & 98.26\\
4 & 6 & 13 & 15 & 513.76 & 11 & 22 & 23 & 346.54\\
8 & 12 & 25 & 30 & 578.3 & 12 & 22 & 28 & 566.48\\
16 & 19 & 29 & 36 & 853.48 & 21 & 31 & 36 & 745.74\\
32 & 28 & 46 & 57 & 1074.35 & 42 & 60 & 76 & 745.67\\
\bottomrule
\end{tabular}
\caption{GetEvidence (query): latency and committed throughput by concurrency (failed transactions: bft4-r2 0, raft4-r2 0).}
\label{tab:bench-getevidencequery}
\end{table}

\begin{table}[htbp]\centering\small
\begin{tabular}{@{}lrrrrrr@{}}
\toprule
Run & Load & \multicolumn{2}{c}{RecordEvent(ACCESS)} & \multicolumn{2}{c}{RegisterEvidence} & Query\\
\cmidrule(lr){3-4}\cmidrule(lr){5-6}\cmidrule(lr){7-7}
 & avg & p50, c=1 & tx/s, c=32 & p50, c=1 & tx/s, c=32 & p50, c=1\\
\midrule
bft4-r2 & 2.46 & 64 ms & 198.35 & 65 ms & 182.11 & 6 ms\\
bft4 & -- & 103 ms & 179.57 & 121 ms & 167.38 & 8 ms\\
raft4-r2 & 2.43 & 105 ms & 195.61 & 105 ms & 131.45 & 11 ms\\
raft4 & -- & 128 ms & 94.81 & 127 ms & 82.55 & 13 ms\\
\bottomrule
\end{tabular}
\caption{All benchmark runs. Round 1 (\code{bft4}, \code{raft4}) ran directly after network builds; round 2 (\code{-r2}) after a 60~s settling period. Queries do not involve ordering, so their variation between runs measures host noise.}
\label{tab:bench-runs}
\end{table}

\paragraph{Results.} Tables \ref{tab:bench-recordeventaccess}--\ref{tab:bench-getevidencequery} give the settled round (round 2); \cref{tab:bench-runs} lists all four runs. No transaction failed in any run.
\begin{itemize}
  \item \textbf{Latency is interactive.} A custody write (access record or registration, including endorsement by two agencies, BFT ordering by four consenters and commit) completes in about 65--180~ms at the median and stays below 0.3~s at the 95th percentile up to 32 concurrent clients. For evidence handling, where a few transactions are submitted per item per day, this is far below any operational limit.
  \item \textbf{Throughput has headroom.} On one laptop the BFT network committed about 180--200 transactions per second at 32 concurrent clients and had not saturated; queries reached about 1{,}000 per second. Evidence handling needs orders of magnitude less.
  \item \textbf{No measurable cost of BFT at this scale.} In the settled round, BFT and Raft are within each other's run-to-run variation (BFT was faster at concurrency 1 and for registrations at 32; access-log throughput at 32 was equal). The first round, run directly after each network build, was slower for both and much slower for Raft; queries, which do not involve ordering at all, varied by about a factor of two between runs. On a single shared host, run-to-run noise is therefore as large as any protocol difference, and the honest conclusion is that \emph{the Byzantine fault tolerance of four agency-operated consenters costs no measurable latency or throughput here}, not that BFT is faster.
  \item \textbf{What would change on a real deployment.} With consenters at four agencies, every ordering round crosses the network; BFT needs more message rounds than Raft, so wide-area latency will add more to BFT than to Raft. Repeating this benchmark across hosts is part of the deployment work (\cref{sec:limits}).
\end{itemize}

\subsection{Threats to validity}
\begin{itemize}[nosep]
  \item \textbf{Single host.} All nodes share one laptop's CPU and disk and there is no network latency between agencies; absolute latencies are optimistic and throughput is bounded by the host, not by the protocol.
  \item \textbf{Synthetic evidence.} Identity-check thresholds come from synthetic captures (\cref{sec:limits}); they demonstrate the mechanism, not its error rates.
  \item \textbf{Workload.} Real evidence handling is low-volume (registrations and transfers per case, not per second); the benchmarks show headroom rather than a target load.
  \item \textbf{Adversary model of the tamper suite.} Attacks were run with the attacker's access simulated (direct file writes, chosen endorsers); compromise of a majority of agencies or of signing keys is out of scope (\cref{sec:limits}).
\end{itemize}
